% Laget automatisk av elevark/lag_elevark.py fra sim/sirkelbevegelse.html. Endringer her blir overskrevet.
\documentclass[11pt]{article}
\usepackage{elevark}
\begin{document}
\elevarkhode{Fysikk 2 · Sirkelbevegelse}{Sirkelbevegelse}{https://fysikklab.no/sim/sirkelbevegelse.html}{Simuleringen}

En ball svinges rundt i en snor med konstant fart. Farten er like stor hele tiden, men retningen endrer seg, og da har ballen akselerasjon. Hvilken vei peker akselerasjonen, og hvor stor er den?

\section*{Oppgaver}

Start med $r = 2{,}0\,\mathrm{m}$, $v = 3{,}0\,\mathrm{m/s}$ og $m = 1{,}0\,\mathrm{kg}$.

\begin{oppgaveliste}
\oppgave \fase{gjett}{Gjett}Farten er konstant. Har ballen likevel akselerasjon? Hvilken vei peker den? Slå på «Vis fartsendringen $Δv$» og se hvilken vei $Δv$ peker.
\svar
\oppgave \fase{test}{Regn}Regn ut omløpstiden, vinkelfarten og sentripetalakselerasjonen. Hvor stor kraft må snora dra med? Sjekk.
\svar
\oppgave \fase{gjett}{Gjett}Hva skjer med akselerasjonen hvis du dobler farten? Hva hvis du i stedet dobler radien og holder farten fast? Og hvis du dobler radien og holder omløpstiden fast?
\svar
\oppgave \fase{forklar}{Forklar}Trykk på «Kutt snora». Hvilken vei går ballen? Forklar med Newtons 1. lov. Hvorfor flyr den ikke rett utover fra sentrum?
\svar
\oppgave \fase{test}{Regn}En bil kjører gjennom en rundkjøring med radius 15~m i 36~km/h. Hva er sentripetalakselerasjonen? Hvor stort må friksjonstallet minst være for at bilen ikke skal skli?
\svar
\oppgave \fase{forklar}{Forklar}Det finnes ingen egen «sentripetalkraft». Hvilken kraft gir sentripetalakselerasjonen for ballen i snora, for bilen i rundkjøringen og for Månen rundt Jorda?
\svar
\oppgave \fase{utfordring}{Utfordring}Månen går rundt Jorda med radius $3{,}84 \cdot 10^{8}\,\mathrm{m}$ og omløpstid 27,3 døgn. Regn ut sentripetalakselerasjonen til Månen. Sammenlign med $g/3600$. Månen er omtrent 60 jordradier unna. Hva sier det om tyngdekraften?
\svar
\end{oppgaveliste}

\end{document}
